\documentclass[11pt]{article}
\usepackage[margin=0.85in]{geometry}
\usepackage{amsmath,amssymb,bm,booktabs,longtable,array,xcolor,hyperref}
\usepackage[T1]{fontenc}
\usepackage{lmodern}
\usepackage{needspace}
\hypersetup{colorlinks=true,linkcolor=blue,urlcolor=blue}
\newcommand{\Tr}{\operatorname{Tr}}
\newcommand{\Rea}{\operatorname{Re}}
\newcommand{\Herm}{\operatorname{Herm}}
\newcommand{\dd}{\partial_\nu}
\newcommand{\Sig}{\widetilde\Sigma}
\newcommand{\Sym}{\mathcal S}
\setlength{\parskip}{5pt}
\setlength{\parindent}{0pt}
\setlength{\emergencystretch}{3em}
\title{The genuinely eDMFT projector contribution to stress\\
\large Derivation, source audit, and independent tests}
\author{Working audit for Kristjan Haule}
\date{November 7, 2026}
\begin{document}
\maketitle
\textbf{Scope.} This audit examines projector 5 and the current sources in
\path{/Users/haule/projects/eDMFT_new/src/dmft2}. It does not repeat the
Blaha G1/G2 derivations. Production sources and original material calculations
were not changed. All new runs used isolated copies.

\textbf{Main findings.} The direct Matsubara gradient, matrix normalization
derivative, and complementary-band response pass independent algebraic and
finite-functional tests. One chain-rule inconsistency is established: the
normalization overlap is symmetrized, but the adjoint symmetry operation in
the normalization-window response is disabled by default. Applying the existing
diagnostic switch changes the SrVO$_3$ mean stress by $+0.1293$ GPa and the
FeSe mean stress by $+0.1397$ GPa; it
does not resolve the hydrostatic discrepancy. Two FeSe shear checks fail and
require a consistent equivalent-atom perturbation convention. A separate,
more fundamental issue is potential dependence of a finite-band projector:
ordinary charge self-consistency need not make that implemented functional
stationary. An independent finite model demonstrates this effect; its size in
the materials has not been measured.

\section{Definitions and matrix orientations}
The derivative $\partial_\nu$ refers to one homogeneous strain channel.
The six channels are $xx,yy,zz,yz,xz,xy$. For a shear channel $ab$, a scalar
perturbation $t$ means $\epsilon_{ab}=\epsilon_{ba}=t/2$. Thus the derivative
with respect to $t$ is $\Omega\sigma_{ab}$, without an extra factor two.
For isotropic length strain, $d\Gamma/dt=\Omega\Tr\sigma$ and the mean stress
is $(3\Omega)^{-1}d\Gamma/dt$. Here the reported sign is $+\partial\Gamma/\Omega$.

\begin{longtable}{p{0.20\linewidth}p{0.74\linewidth}}
\toprule Symbol & Definition\\\midrule
$k,w_k$ & Reduced wave vector and normalized Brillouin-zone weight;
$\partial_\nu w_k=0$. Spin is an explicit index or a degeneracy factor in
the functional sums, but not an extra factor in the overlap normalization.\\
$\chi_{Ik}$ & LAPW basis function; $I,J$ label basis functions.\\
$O_k,H_k$ & Basis overlap and KS Hamiltonian;
$H_k A_k=O_k A_k\varepsilon_k$, $A_k^\dagger O_k A_k=1$.\\
$n,m,a$ & KS band indices. $n,m\in W_k$ when specified; $a\notin W_k$
denotes a complementary band.\\
$\alpha,\beta$ & Local correlated orbitals. Independent correlated blocks
are summed; formulas below display one block for clarity.\\
$\phi_\alpha$ & Chosen localized orbital, with its radial shape and local-axis
convention specified.\\
$q_{I\alpha}$ & $\langle\chi_I|\phi_\alpha\rangle$, a basis-by-orbital matrix.\\
$P_k$ & Raw band-by-orbital projector $A_{Wk}^\dagger q_k$.\\
$S,X,U_k$ & Global overlap, $X=S^{-1/2}$, and normalized projector $U_k=P_kX$.\\
$\Sig$ & $\Sigma_{\rm imp}-V_{\rm DC}$ in the correlated orbital basis.\\
$G_k(z)$ & Band Green function, $[z-\varepsilon_k-U_k\Sig(z)U_k^\dagger]^{-1}$,
with the self-energy embedded only in $W_k$; $z=i\omega_n+\mu$.\\
$\rho_k,\rho^E_k$ & Density and energy-weighted density matrices with the same
frequency summation, spin and $k$ weights. $\rho^E$ uses the regularized
moment $zG-1$. Both matrices are Hermitian after pairing frequencies.\\
$C_k$ & Orbital-by-window-band functional gradient, defined in Eq.~\eqref{eq:C}.
It is not a projector and is not generally Hermitian.\\
$\Sym$ & The linear symmetry reconstruction/averaging applied to the
irreducible-$k$ overlap. $\Sym^*$ is its adjoint under the real trace pairing.\\
\bottomrule
\end{longtable}
All formulas assume a positive overlap on the active correlated subspace and
a spectrally separated boundary of the fixed-index window. Neither condition
requires $U_k$ to have full column rank at every individual $k$.

\section{Start from the functional, not the total energy}
Write the stationary functional using independent variables $n,V,g,\Sig$:
\begin{align}
\Gamma={}&F_{\rm sp}[H_0+V,\Sig,U]+E_{\rm Hxc}[n]-\int Vn
 +E_{\rm ion}\notag\\[-2pt]
 &+\Phi_{\rm imp}[g]-\Phi_{\rm DC}[g]-\Tr(\Sig g).
\label{eq:Gamma}
\end{align}
Here $H_0$ includes the kinetic operator and ionic external potential, $V$ is
the Hartree--XC potential, and $F_{\rm sp}$ is the regularized lattice
$-T\sum\ln\det[-G^{-1}]$, including $\mu N$ for a fixed-particle-number
Helmholtz functional. The local trace includes frequency, block and spin sums.
This is a functional before imposing stationarity, not a rearrangement of
the nonstationary eDMFT total energy.

At the stationary impurity solution,
\begin{equation}
\frac{\delta(\Phi_{\rm imp}-\Phi_{\rm DC})}{\delta g}=\Sig,
\qquad g=\sum_k w_k U_k^\dagger G_k U_k.
\end{equation}
Consequently the $\Sig\,\delta g$ and $g\,\delta\Sig$ variations cancel
between the local functional and the lattice term. A two-particle vertex
is not needed for this first derivative. Explicit variation of $U$ survives:
\begin{align}
\delta\Sigma_{\rm emb}
 &=\delta U\Sig U^\dagger+U\Sig\delta U^\dagger+U\delta\Sig U^\dagger,\\
D^{\rm proj}_\nu
 &=2\Rea\sum_k\Tr[C_k\,\dd U_k],\label{eq:D}\\
 C_k&=w_k T\sum_{\omega_n}\Sig(i\omega_n)U_k^\dagger G_k(i\omega_n).
\label{eq:C}
\end{align}
Spin sums or their degeneracy factors are included in $C$. The $k$ weight
must not be applied a second time. Negative Matsubara frequencies supply the
conjugate partner needed for the real expression. A fixed self-energy is
therefore sufficient to test this term, with the impurity functional replaced
by its tangent plane. This leaves $F_{\rm sp}$ plus the density/ionic terms
and a strain-independent constant; it does not eliminate the projector chain.

In a moving nonorthogonal basis the full one-particle variation is
\begin{equation}
D^{\rm 1p}_\nu=\sum_k\Rea\Tr[\rho_k H_{k,\nu}-\rho^E_k O_{k,\nu}]
                  +D^{\rm proj,explicit}_\nu .\label{eq:master}
\end{equation}
The meaning of ``explicit'' must match the representation: once KS-state
connections have been absorbed into the matrix pair, those same connections
cannot be added a second time through $\dd U$.

\section{The direct gradient and the $B=UC$ identity}
Dyson's equation gives, with a consistent regularization,
\begin{equation}
B_k\equiv\rho^E_k-\varepsilon_k\rho_k=U_k C_k.\label{eq:B}
\end{equation}
The KS energy matrix multiplies $\rho$ on the \emph{left}. $B$ need not be
Hermitian. Recovering $C$ by a pseudoinverse of $U$ can lose its component
in $\ker U$. The current production code instead accumulates $C$ directly.
For a positive Matsubara frequency its paired integrand is
\begin{equation}
C^{\rm pair}(i\omega)=
\Sig U^\dagger G+\Sig^\dagger U^\dagger G^\dagger.\label{eq:pair}
\end{equation}
The second term is not the adjoint of the first rectangular matrix.
\texttt{AccumulateDirectProjectorPair} implements these products in this order.
Its static-reference and high-frequency corrections are assembled with the
same conventions as $\rho^E$. Tests differentiate a complex log determinant
and include a rank-deficient projector; both pass. Actual SrVO$_3$ and FeSe
diagnostics also compare $B$ against an independently accumulated
$\Sigma_{\rm emb}G$ moment.

\section{Differentiate the actual matrix normalization}
The implemented overlap and normalized projector are
\begin{equation}
S=\Sym\!\left[\sum_k w_k P_k^\dagger P_k\right],\qquad
X=S^{-1/2},\qquad U_k=P_kX.\label{eq:norm}
\end{equation}
Thus
\begin{align}
\dd U_k&=(\dd P_k)X+P_k(\dd X),\\
\dd S&=\Sym\!\left[\sum_k w_k
 ((\dd P_k)^\dagger P_k+P_k^\dagger\dd P_k)\right].\label{eq:dS}
\end{align}
For a symmetry-preserving strain the symmetry map is fixed. For general
strain the Cartesian derivative indices must be transformed as well; one
must use the adjoint of that complete map, not assume independent scalar
symmetrization of each shear component.

Put $Y=X^{-1}$. Differentiating $XSX=1$ gives the Sylvester equation
\begin{equation}
Y(\dd X)+(\dd X)Y=-X(\dd S)X.\label{eq:sylvester}
\end{equation}
If $X=V_x\operatorname{diag}(x_i)V_x^\dagger$, primes on matrices in the
next equation mean transformation to this eigenbasis:
\begin{equation}
(\dd X)'_{ij}=-\frac{x_i x_j(\dd S)'_{ij}}{1/x_i+1/x_j},\qquad
(X^{-1}\dd X)'_{ij}=-\frac{x_j(\dd S)'_{ij}}{1/x_i+1/x_j}.
\end{equation}
No derivative of the eigenvectors of $X$ is required. In particular,
$\dd X=-\tfrac12S^{-3/2}\dd S$ is invalid unless the matrices commute.

An economical way to carry the normalization response is its adjoint.
Define the Hermitian orbital matrix and the linear Sylvester map
\begin{equation}
J=\Herm\sum_k C_kP_k,\qquad \mathcal L_X[\delta S]=\delta X,
\qquad Q=\mathcal L_X[J].
\end{equation}
$\mathcal L_X$ is self-adjoint on Hermitian matrices. Therefore
\begin{align}
D^{X}_\nu&=2\Tr[J\dd X]=2\Tr[Q\dd S]\\
 &=2\Tr\!\left[\underbrace{\Sym^*[Q]}_{Q_{\rm eff}}
       \sum_k w_k\big((\dd P_k)^\dagger P_k+P_k^\dagger\dd P_k\big)\right].
\label{eq:adjoint}
\end{align}
An especially useful complete gradient with respect to the raw projector is
\begin{equation}
\boxed{\mathcal C_{P,k}=XC_k+2w_kQ_{\rm eff}P_k^\dagger,\qquad
D^{\rm proj}_\nu=2\Rea\sum_k\Tr[\mathcal C_{P,k}\dd P_k].}\label{eq:rawgrad}
\end{equation}
The source stores the equivalent moment $(\sum_k C_kU_k)^T$ and uses
$X^{-1}\dd X$. Its conjugations agree with the above formula; comments using
a band-by-orbital Frobenius gradient use a different orientation of $C$.

\textbf{Confirmed chain-rule omission.} The forward overlap is symmetrized,
while \texttt{Symmetrize\allowbreak{}Normalization\allowbreak{}Moment} is default-off. In the
symmetry-invariant case, projecting the moment before the Sylvester adjoint
is equivalent to Eq.~\eqref{eq:adjoint}. In general the robust construction
is to apply the exact adjoint after it, respecting active blocks and the
order of equivalent-atom averaging. The diagnostic switch
\texttt{EDMFT\_STRESS\_NORM\_SYMMETRY=APPLY} tests the existing implementation;
changing that default alone is not a proof of all general shear conventions.

\section{Fixed band indices still have a window derivative}
Let $A_n$ be a normalized generalized KS eigenvector. In the reference band
basis define $T_{an}=A_a^\dagger O\dd A_n$. For separated energies,
\begin{equation}
T_{an}=\frac{H_{\nu,an}-\varepsilon_n O_{\nu,an}}
                 {\varepsilon_n-\varepsilon_a}\quad(a\ne n),\qquad
T_{nn}=-\tfrac12O_{\nu,nn},\label{eq:T}
\end{equation}
where the diagonal imaginary gauge is chosen to vanish. Then
\begin{equation}
\dd P_W=T_{F,W}^\dagger P_F+A_W^\dagger\dd q.\label{eq:dP}
\end{equation}
$F$ is the full available band set, not just the window. There is no moving
energy-cutoff delta function for fixed band indices. Nevertheless, the
selected spectral subspace moves, and its complementary bands enter
Eq.~\eqref{eq:T}. At a degeneracy crossing a window edge, the spectral
projector is not differentiable without a prescription that includes the
complete degenerate subspace. Silently skipping a small denominator is not
a physical resolution of that problem.

At fixed $X$, put $K_{an}=(U_F C)_{an}$. The outside-window contribution is
\begin{equation}
D^{W}_\nu=2\Rea\sum_k\sum_{n\in W_k,a\notin W_k}
  \frac{K_{an}\,[H_{\nu,an}-\varepsilon_nO_{\nu,an}]^*}
       {\varepsilon_n-\varepsilon_a}.\label{eq:window}
\end{equation}
It is represented by Hermitian response matrices
\begin{align}
R_{an}&=K_{an}/(\varepsilon_n-\varepsilon_a),&R_{na}&=R_{an}^*,\\
R^E_{an}&=\varepsilon_nR_{an},&R^E_{na}&=(R^E_{an})^*,\\
D^{W}_\nu&=\Rea\sum_k\Tr[R_k H_{k,\nu}-R^E_k O_{k,\nu}].
\end{align}
Putting $K_{an}$ in the transposed matrix slot breaks complex KS-phase
invariance. The current source has the correct orientation and passes a
rephasing test.

The normalization window response uses the same construction with the raw
projector and $C^{N}_{P,k}=2w_kQ_{\rm eff}P_k^\dagger$. It also needs the
within-window metric term
\begin{equation}
\dd S_{\rm metric}=-\Sym\sum_k w_kP_W^\dagger O_{\nu,WW}P_W.
\label{eq:normmetric}
\end{equation}
The source includes it through $R^{E,N}_{WW}=\Herm(P_W C^N_{P,k})$.
Internal unitary rotations cancel from $P_W^\dagger P_W$; the metric term
does not. The remaining within-window connection is already contained in
Eq.~\eqref{eq:master}, using $B=UC$. Adding it again as a separate correction
would double count it.

\section{Raw projector 5, radial response, and the sphere boundary}
Schematically the local raw projector is
\begin{equation}
P_{n\alpha}=\sum_{\lambda}a_{\lambda n}^* I_{\lambda\alpha}
                 +P^{\rm interstitial}_{n\alpha},\qquad
I_{\lambda\alpha}=\int_0^{R_{\rm MT}}r^2
 u_\lambda(r)^*\phi_\alpha(r)\,dr .
\end{equation}
Angular, spin and local-axis transformations are understood. The primitive
index $\lambda$ includes $u_l,\dot u_l$ and all local orbitals. Consequently
\begin{equation}
\dd P_{n\alpha}=
\sum_\lambda(\dd a_{\lambda n})^* I_{\lambda\alpha}
 +\sum_\lambda a_{\lambda n}^*\dd I_{\lambda\alpha}
 +\dd P^{\rm interstitial}_{n\alpha}.\label{eq:radial}
\end{equation}
The production builder includes matching/Bessel changes, angular derivatives,
the $-P/2$ volume term in each diagonal channel, and the projector-5
interstitial continuation when present. It also has a branch for
$\dd I$, but \texttt{drix\_mat} is zero in the normal production path.
That is a frozen-radial approximation, not proof that $\dd I=0$ along a
charge-relaxed strain path. Its repair must be paired with the same radial
response in $H_\nu$ and $O_\nu$.

A projector finite at $R_{\rm MT}$ does not automatically imply a new
independent surface correction. With fixed physical radius and coordinates
relative to the co-moving atom, the integration domain is fixed. If the
radius changes, there is explicitly
\begin{equation}
(\dd I)_{\rm boundary}=R_{\rm MT}^2(\dd R_{\rm MT})
 u_\lambda(R_{\rm MT})^*\phi_\alpha(R_{\rm MT}),
\end{equation}
with the angular contraction included as appropriate. A laboratory-coordinate
derivation must combine moving-boundary and translation terms consistently.
For a projector continued into the interstitial, the inside/outside boundary
pieces must be combined before claiming an additional surface term.

\section{A stationarity issue specific to spectral projectors}
This section is a separate derivation, not a claim that the observed GPa
error has already been explained. In a fixed Hilbert space define
\begin{equation}
\Pi_W[V]=\sum_{n\in W}|\psi_n[V]\rangle\langle\psi_n[V]|,
\qquad |\Phi_W[V]\rangle=\Pi_W[V]|\phi\rangle X[V].
\end{equation}
The finite-window embedding is a physical function of $V$, even when the
radial orbital on disk is fixed. A basis change alone is different and
cancels exactly between the terms in Eq.~\eqref{eq:master}.

Varying Eq.~\eqref{eq:Gamma} at fixed geometry gives
\begin{equation}
\delta_V\Gamma=\int(n_G-n)\delta V
       +2\Rea\sum_k\Tr[C_k\delta_V U_k]_{\rm physical\ window}.
\label{eq:Vstationary}
\end{equation}
Ordinary charge self-consistency sets $n=n_G$. It does not set the second
term to zero. At fixed basis, $\delta O=0$, $\delta q=0$, this residual is
\begin{equation}
\delta_V\Gamma=\Rea\sum_k\Tr[(R_k+R_k^N)\delta V_k].\label{eq:RV}
\end{equation}
The matrices are exactly the complementary-window and normalization-window
responses above. The inspected code uses them in force/stress contractions,
but constructs the charge from the ordinary lattice density matrix.

It follows that a derivative along a converged charge path may contain
\begin{equation}
\boxed{D^{\rm path}_\nu
 =D^{\rm explicit}_\nu+
 \Rea\sum_k\Tr[(R_k+R_k^N)(\partial_\nu V_{\rm Hxc})_k]
 +\text{other projector/basis response terms}.}\label{eq:path}
\end{equation}
Here the extra $V$ derivative means the induced potential variation excluded
from the explicit strain convention, not a second copy of a potential
derivative already present in $H_\nu$. The exact split must be verified
against the LAPW stress kernels. Complete stationary variation would instead
require the projector contribution in the $V$ stationarity equation, or a
projector definition genuinely independent of $V$. Either choice changes
the formulation and must not be made silently as a stress-only patch.

An independent dynamic-self-energy finite model verifies
Eq.~\eqref{eq:Vstationary} without radial approximations or core/valence
cancellation:
\begin{center}
\begin{tabular}{lrr}\toprule
Test & Finite $W$ & Complete space\\\midrule
$\delta F_{\rm sp}-\Tr(\rho\delta V)$ & $0.035133243630$ & $-6.5\,10^{-11}$\\
Predicted window chain & $0.035133243587$ & $0$\\\bottomrule
\end{tabular}
\end{center}
These are model energy-derivative units, not GPa. The density equals the
exact lattice density at the reference point. Thus the residual is not
an SCF tolerance error. The material value of Eq.~\eqref{eq:path} remains
unmeasured, and a full-space model does not establish completeness of a
finite LAPW basis.

\Needspace{7\baselineskip}
\section{Tests and implementation status}
The regression harness extracts production Fortran routines and compiles
them with bounds checking. Independent Python endpoints diagonalize the
generalized Hamiltonian and normalization matrices. The joint test uses
three unequal $k$ weights, complex projectors, both lower and upper
complementary bands, nonzero $H_\nu,O_\nu,q_\nu$, and noncommuting matrices.
A second version uses an exactly soluble dynamical self-energy
$\Sig(z)=\Sigma_\infty+V_b(z-E_b)^{-1}V_b^\dagger$; its free energy is obtained
from the enlarged bath Hamiltonian without a Matsubara-tail approximation.

\begin{center}\small
\begin{tabular}{p{.60\linewidth}p{.30\linewidth}}\toprule
Test & Result\\\midrule
Inverse square root against endpoint diagonalization & Pass, $4.7\,10^{-12}$\\
Normalization adjoint & Pass, $1.6\,10^{-11}$\\
Complex window response and KS rephasing & Pass\\
Generalized $H_\nu/O_\nu$ window contraction & Pass, $1.5\,10^{-11}$\\
Within-window normalization metric & Pass, $3.4\,10^{-11}$\\
Pure basis-connection cancellation & Pass, $3.5\,10^{-17}$\\
Direct $C$, paired Matsubara / real axis & Pass, below $3\,10^{-11}$\\
Rank-deficient $U$ with direct $C$ & Pass, $2.1\,10^{-11}$\\
All projector chains, no symmetry map & Pass, below $10^{-10}$\\
All chains, symmetrized $S$, default response & Fail\\
Same test with adjoint symmetry & Pass, below $10^{-10}$\\
Potential stationarity with a finite window & Nonzero residual reproduced\\
Same stationarity test with the complete space & Residual below $10^{-10}$\\\bottomrule
\end{tabular}\end{center}

For example, in the joint dynamical test with symmetrized $S$, the finite
derivative is $-0.039320865142$. The default chain gives $-0.099349880276$;
the adjoint-complete chain gives $-0.039320865133$. This proves a chain-rule
failure for the tested map, without fitting a stress correction to a material.
The failure persists with a symmetry-preserving dynamical self-energy:
finite difference $-0.105243690363$, default chain $-0.105612238523$,
adjoint-complete chain $-0.105243690327$. It is therefore not an artifact
of testing a symmetry-breaking self-energy against a symmetrized overlap.

\subsection*{Actual material checks}
\begin{center}
\begin{tabular}{lrr}\toprule
Check & SrVO$_3$ & FeSe RELA\\\midrule
Maximum diagonal raw-projector relative error & $5.81\,10^{-10}$ & $1.15\,10^{-9}$\\
Maximum shear raw-projector relative error & $7.59\,10^{-10}$ & $0.5802$\\
Symmetry-adjoint mean-stress change (GPa) & $+0.1293$ & $+0.1397$\\\bottomrule
\end{tabular}
\end{center}
The switch changes the normalization-window response. Neither the impurity
self-energy nor the charge inputs were changed. SUM, VSUM, YSUM and DRHO are
identical between the two modes in each material. The SrVO$_3$ hydrostatic
mismatch remains; the switch slightly increases it. These are changes in
\emph{mean} stress, already divided by three, with the final mixer sign.


The independently accumulated $\Sigma_{\rm emb}G$ moment agrees with
$\rho^E-\varepsilon\rho$ to maximum relative errors $6.82\,10^{-14}$ in
SrVO$_3$ and $3.46\,10^{-14}$ in FeSe. Direct $C$ and its pseudoinverse
reconstruction agree to $2.22\,10^{-14}$ and $1.19\,10^{-14}$, respectively,
in these runs. Thus a lost pseudoinverse component is not the current
material discrepancy. A separate rank-deficient model still requires direct
$C$, and the production code correctly retains it.

Reusing the saved central density and core contributions, the fresh valence
runs give mean stresses $5.54758\to5.67688$ GPa in SrVO$_3$ and
$-0.25286\to-0.11317$ GPa in FeSe when the switch is applied. The previously
audited numerical references are about $2.934$ and $-1.0491$ GPa,
respectively. The symmetry step therefore slightly worsens both comparisons.
These small source/assembly checks are not newly reconverged strain triplets.

All four full material runs completed and the audited outputs were finite.
Their Fortran termination messages included sticky invalid/divide-by-zero
IEEE flags; the origin of those flags was not traced in this pass.

The raw-projector checks hold eigenvector coefficients and radial functions
fixed, use the first irreducible $k$ point, and perturb geometry by $10^{-6}$.
They test the explicit builder, not the complete relaxed projector. For FeSe,
the failed $xz$ and $xy$ channels are exactly those reversed by the operation
$\operatorname{diag}(1,-1,-1)$ relating the two Fe atoms. The builder's
matching magnitudes use shared Bessel arrays computed from $|B(k+G)|$,
whereas some analytic factors use $|B R(k+G)|$. These agree at the symmetric
reference and under diagonal strain, but differ under the symmetry-breaking
shears. The finite-difference helper retains the original equivalence map.
This identifies a convention inconsistency to isolate; it does not by itself
decide whether the physical shear formula or the restricted helper should be
changed. It is not an explanation of the hydrostatic discrepancy.

Separate scratch calculations retaining only the first or only the second
correlated Fe atom confirm the localization. Atom 1 passes every channel
with maximum relative error $1.39\,10^{-9}$. Atom 2 has relative errors
$0.7148$ and $0.7805$ in $xz$ and $xy$, respectively; its other four channels
pass at about $10^{-9}$. The absolute discrepancy norms are exactly the
same as in the two-atom test. These one-$k$ probes test projector geometry
only and are not meaningful material pressure calculations.

\Needspace{10\baselineskip}
\section{Source map and the next controlled tests}
\begingroup\small
\begin{longtable}{p{.46\linewidth}p{.48\linewidth}}
\toprule Routine / file & Audited role\\\midrule
\texttt{cmp\_dmft\_weights.f90} & Direct $C$, positive/negative-frequency pairing,
energy moments, static reference and tails.\\
\texttt{Build\_DMFT\_Projector} in \texttt{l2mainSubs.f90} & Explicit raw projector,
matching/angular/volume terms; optional $\partial I$.\\
\texttt{cmp\_overlap} & $k$-summed raw overlap and explicit derivative,
symmetry reconstruction, $X$, preservation of a saved normalization.\\
\texttt{ComputeNormalizationDerivative} & Sylvester derivative, Eq.~\eqref{eq:sylvester}.\\
\texttt{BuildWindowResponseMatrices} & Direct $C$ selection and fixed-$X$
complementary-band matrices.\\
\texttt{AccumulateNormalizationMoment}, \texttt{BuildNormalizationAdjoint} &
Global $k$-coupled normalization gradient.\\
\texttt{BuildNormalizationWindowResponseMatrices} & Complementary normalization
response plus the within-window metric.\\
\texttt{symmetrize.f90} & Forward overlap symmetry; adjoint-moment switch defaults
to off; Cartesian derivative rotations and equivalent-block averaging.\\
\texttt{l2main.f90} & Calls the above, contracts $2\Rea C\partial U$, and applies
the negative sign of the intermediate WIEN \texttt{STR\_VAL} convention.\\
\texttt{p\_project.f90} & Projector radial overlaps; production
\texttt{drix\_mat=0}.\\\bottomrule
\end{longtable}
\endgroup

\Needspace{7\baselineskip}
\textbf{Recommended next work, requiring approval before source edits:}
\begin{enumerate}
\item Implement the exact adjoint of the actual normalization/symmetry map,
including the order of equivalent-block averaging. Keep a regression test
with noncommuting complex matrices and a dynamical self-energy. Its necessity
is mathematical; the sign must not be chosen to improve a material number.
\item Isolate the FeSe shear check by equivalent atom and separate Bessel,
angular and LO pieces. For broken symmetry use a consistent full-cell
perturbation or rotate the strain together with the atom mapping.
\item At one fixed SrVO$_3$ geometry, perturb the input density/potential in
a charge-neutral direction and evaluate the frozen-self-energy tangent
functional. Compare its slope with $\Tr[(R+R^N)\delta V]$ using the same
potential matrices. Include a DFT control to separate radial-basis effects.
This directly tests the stationarity gap; it avoids subtracting large
independent strained core and valence energies.
\item Only then contract the measured induced $V_{\rm Hxc}$ response along
the strain path. Establish which part is already in the existing kernels
before adding Eq.~\eqref{eq:path} to any reported stress.
\item Complete the paired radial $H/O/P$ response and repeat stored-band
convergence. Stored complementary bands remain an approximation; increasing
their number does not enlarge the chosen DMFT window itself.
\end{enumerate}
Other limitations remain: window-edge degeneracies, local-axis conventions
with a diagonal-only self-energy, an outdated saved normalization, and
consistency of active-orbital selection. The forward normalization uses
nonzero diagonal \texttt{Sigind} entries, while some response routines accept
any active row/column; a purely off-diagonal active block is therefore an
additional unsupported corner case. It is not the present SrVO$_3$/FeSe setup.

An adjacent force-path limitation also deserves a separate audit: the
global normalization-window pass is gated by \texttt{Qstress}, whereas the
fixed-$X$ window response is also used for forces. Explicit force
normalization derivatives exist, but this does not establish that the full
normalization-window force chain is present. No new force benchmark was
performed in this pass.

\textbf{Reproducibility.} The accompanying Python scripts and JSON results
record the combined functional test, potential-stationarity test, material
diagnostics, source hashes, and copies of the small diagnostic outputs.
They do not change the production executable or commit any source changes.
\end{document}
